\section{27 сентября 1914 г.}

\lp{20260604_193456434}

Милый, пишу все понемногу, потому что каждый божий день какое-нибудь неотложное
дело или бесконечный жаркий разговор с кем-нибудь. Всем приходится
устанавливать со всеми известное отношение: то, что сейчас происходит
заставляет людей выяснять очень многое и между собой. Вчера вечером был у нас
Николай Петрович. Прочла ему кое-что из ваших писем (что касается Юнга, то не
беспокойтесь: никто (кроме Ильина) не услышит от меня о вашем лечении). Он был
очень благодарен за это и во многом со всем согласен с вами. Но я очень вас
прошу не писать никому о событиях и мне тоже, а то письма не дойдут. И вообще и
по другим соображениям не советую вам пока высказываться ни с кем. Книгу К. вам
выслал через Христианию. Ее оттуда перешлют вам под новой бандеролью, а то тут
не принимали в Швейцарию. О делах Мусагета пока не беспокойтесь. Поговорю с
Марго, но у нее теперь огромные расходы. К. говорит, что пока еще это не так
необходимо. С Колей я говорила, пока еще он не может выплачивать, но как
только

\lp{20260604_193435130}

\noindent
явится возможность, он будет погашать на первом плане именно этот долг.

\lp{20260604_193445066}

Когда я уезжала в Крым, то просила вас отдать Мане 20 р.
Сделали вы это или нет?
Я не помню, говорили ли вы мне об этом.
